\section{Le endomembrane}

Il sistema delle \textbf{endomembrane} comprende: reticolo endoplasmatico, complesso di Golgi, lisosomi, perossisomi, vacuoli.

% ---------------------------------------------------------------------------
\subsection{Reticolo endoplasmatico: struttura}

\rubrica{Liscio (REL)}
\begin{itemize}
  \item costituito da membrane tubulari;
  \item sistema di canali interconnessi che si diramano nel citoplasma.
\end{itemize}

\rubrica{Rugoso (RER)}
\begin{itemize}
  \item costituito da sacche appiattite (\textit{cisterne});
  \item presenza di ribosomi sulla superficie esterna;
  \item è continuo con la membrana esterna dell'involucro nucleare.
\end{itemize}

% ---------------------------------------------------------------------------
\subsection{Reticolo endoplasmatico liscio: funzioni}

\begin{itemize}
  \item biogenesi di membrane;
  \item sintesi di colesterolo e ormoni steroidei;
  \item nel fegato: detossificazione di composti organici;
  \item sequestro di ioni calcio nelle cellule muscolari.
\end{itemize}

% ---------------------------------------------------------------------------
\subsection{Reticolo endoplasmatico rugoso: funzioni}

\begin{itemize}
  \item \textbf{sintesi delle proteine}: secrete dalla cellula, integrali di membrana, solubili situate all'interno di organuli cellulari;
  \item sintesi delle catene di carboidrati;
  \item sintesi dei fosfolipidi.
\end{itemize}

\rubrica{Sintesi di una proteina di secrezione o di un enzima lisosomale}
\begin{enumerate}
  \item la sintesi inizia su un ribosoma libero; quando la sequenza segnale emerge dal ribosoma si lega alla \textbf{SRP} (particella di riconoscimento del segnale), che blocca la traduzione;
  \item il complesso catena nascente--ribosoma--SRP prende contatto con la membrana del RE mediante un recettore per SRP;
  \item rilascio dell'SRP e associazione del ribosoma con il \textbf{traslocone};
  \item il peptide segnale si lega all'interno del traslocone, spostando il tappo del canale e permettendo al polipeptide di attraversare la membrana; quando il polipeptide nascente è entrato nel lume del RE, il peptide segnale viene tagliato da una proteina di membrana e la proteina si ripiega in modo corretto con l'aiuto di proteine specifiche.
\end{enumerate}

\rubrica{Sintesi di una proteina integrale di membrana}
Modello per una proteina con un singolo segmento transmembrana vicino all'estremità N-terminale del polipeptide nascente.
\begin{itemize}
  \item il polipeptide nascente entra nel traslocone come se fosse una proteina di secrezione; l'entrata della sequenza idrofobica transmembrana blocca l'ulteriore traslocazione del polipeptide nascente attraverso il canale;
  \item il traslocone, aprendosi lateralmente, espelle il segmento transmembrana nel doppio strato lipidico;
  \item nella riorientazione il traslocone dispone il segmento transmembrana attraverso le interazioni con le cariche opposte (positive e negative) presenti sulla sua superficie;
  \item il traslocone si apre lateralmente ed espelle il segmento transmembrana nel bilayer, con la disposizione finale della proteina: estremità C-terminale nel lume, estremità N-terminale nel citosol (o viceversa, a seconda dell'orientamento).
\end{itemize}

\rubrica{Mantenimento dell'asimmetria della membrana}
Quando una proteina è sintetizzata nel RER, viene subito inserita nel doppio strato lipidico in un orientamento prevedibile determinato dalla sua sequenza amminoacidica. Questo orientamento è mantenuto durante il suo viaggio nel sistema delle endomembrane. Le catene di carboidrati, aggiunte per prime nel RE, forniscono un modo conveniente per determinare l'asimmetria della membrana, perché sono sempre presenti sul lato cisternale delle membrane citoplasmatiche, che diventa il lato esoplasmico della membrana plasmatica dopo la fusione delle vescicole con la membrana cellulare.

% ---------------------------------------------------------------------------
\subsection{Complesso di Golgi: struttura}

\begin{itemize}
  \item localizzato tra il reticolo endoplasmatico e la membrana plasmatica;
  \item cisterne appiattite discoidali, curve, disposte in pile ordinate e interconnesse da tubuli membranosi;
  \item bordi dilatati associati a vescicole che gemmano dalla parte periferica di ogni cisterna.
\end{itemize}

Organizzazione delle cisterne, da un polo all'altro:
\[
\text{rete } cis \text{ del Golgi (CGN)} \rightarrow \text{cisterne } cis \rightarrow \text{cisterne mediali} \rightarrow \text{cisterne } trans \rightarrow \text{rete } trans \text{ del Golgi (TGN)}.
\]

% ---------------------------------------------------------------------------
\subsection{Complesso di Golgi: funzioni}

\begin{itemize}
  \item \textbf{modificazioni chimiche delle proteine}: rimozione di porzioni di sequenza amminoacidica; aggiunta di gruppi funzionali; aggiunta di unità lipidiche o glucidiche;
  \item \textbf{regolazione del movimento delle proteine}: proteine di secrezione, proteine di membrana, proteine destinate ad altri organuli (es.\ lisosomi).
\end{itemize}

% ---------------------------------------------------------------------------
\subsection{Trasporto delle proteine all'interno della cellula}

\begin{enumerate}
  \item i neopolipeptidi entrano nel RER;
  \item aggiunta di carboidrati a formare glicoproteine;
  \item vescicole di trasporto le trasferiscono sulla superficie \textit{cis} del complesso di Golgi;
  \item ulteriori modificazioni;
  \item le glicoproteine si portano sulla superficie \textit{trans}, dove sono impacchettate su vescicole di trasporto;
  \item vanno sulla membrana plasmatica (o su altri organuli);
  \item il contenuto della vescicola è rilasciato dalla cellula.
\end{enumerate}

% ---------------------------------------------------------------------------
\subsection{Lisosomi}

\begin{itemize}
  \item organelli digestivi;
  \item contengono $\sim$50 enzimi idrolitici: fosfatasi, nucleasi, proteasi, lipasi acida e fosfolipasi, polisaccaridasi e oligosaccaridasi;
  \item attività ottimale a pH acido;
  \item pompa protonica sulla membrana.
\end{itemize}

\rubrica{Funzioni}
\begin{itemize}
  \item \textbf{degradazione delle sostanze} che arrivano dall'ambiente esterno;
  \item \textbf{turnover degli organelli}:
    \begin{enumerate}
      \item un organello viene circondato da una membrana fornita da una cisterna del RE;
      \item si fonde con un lisosoma a formare un \textit{autofagolisosoma};
      \item l'organello viene degradato e i prodotti resi disponibili per la cellula;
      \item al termine del processo l'organello diventa un \textit{corpo residuo};
      \item il corpo residuo può seguire due vie alternative: andare incontro a esocitosi, oppure rimanere nella cellula indefinitamente sotto forma di \textit{granulo di lipofuscina}.
    \end{enumerate}
\end{itemize}

\rubrica{Autofagia}
Un mitocondrio e un perossisoma vengono racchiusi in un involucro a doppia membrana derivato dal RE: si forma così un vacuolo autofagico (autofagolisosoma), che deve fondersi con un lisosoma perché il suo contenuto sia digerito.

% ---------------------------------------------------------------------------
\subsection{Perossisomi}

Detti anche \textit{microcorpi}.
\begin{itemize}
  \item vescicole circondate da membrana;
  \item $\phi$ 0,1--1,0\,$\mu$m;
  \item contengono enzimi ossidativi;
  \item sito di formazione e degradazione di $H_2O_2$;
  \item si duplicano per scissione da organelli preesistenti.
\end{itemize}

\rubrica{Meccanismo redox}
I perossisomi contengono enzimi che riducono l'ossigeno molecolare ad acqua in due tappe:
\begin{itemize}
  \item prima tappa (\textbf{ossidasi}): rimuove elettroni da substrati differenti ($RH_2$, es.\ acido urico o amminoacidi) $\rightarrow$ $R_2 + H_2O_2$;
  \item seconda tappa (\textbf{catalasi}): converte in acqua il perossido di idrogeno formatosi nella prima tappa $\rightarrow$ $2H_2O_2 \rightarrow O_2 + 2H_2O$.
\end{itemize}

% ---------------------------------------------------------------------------
\subsection{Vacuolo delle cellule vegetali}

\begin{itemize}
  \item peculiare delle cellule vegetali;
  \item occupa il 90\% della cellula;
  \item delimitato da una singola membrana: il \textit{tonoplasto};
  \item all'interno alta concentrazione di ioni $\rightarrow$ l'acqua entra per osmosi $\rightarrow$ turgore;
  \item siti di digestione intracellulare;
  \item pH acido per la presenza di una pompa protonica sul tonoplasto.
\end{itemize}
